\section{Related work}
\label{sec:related_work}

Kolmogorov-Arnold netwoks, as interpretable and parameter efficient alternatives of multi-layer perceptrons (MLP), have experienced a great growth in the last year, with the adaptation of existing technologies as convolutional networks \citep{bodner2024convolutional}, residual networks \citep{yu2024residual}, quantum networks \citep{wakaura2025enhanced}, or autoencoders \citep{moradi2024kolmogorovarnoldnetworkautoencoders}. In the same manner, they have had an impressive adoption in domains where transparency is critical, as predicting risks in cardiovascular diseases \citep{al2025kolmogorov} in healthcare \citep{pendyala25effectiveness, pati25kaam}; predicting volatility in finances \citep{CHO2025128781}; predicting biomarkers \citep{alharbi25biomarker}, or predicting physics in power systems \citep{Shuai_2025_powersystems}.

Interpretability is widely regarded as essential for the adoption of ML in sensitive domains \citep{DoshiVelez2017TowardsAR}. Methods for improving interpretability based on KANs have been used in critical domains, such as Kolmogorov-Arnold Additive Models (KAAMs) \citep{pati25kaam}, which are shallow KANs that yields nonlinear additive models, which are well known as interpretable flexible estimators, as Neural Additive Models (NAMs) \citep{agarwal21neural} or generalized additive models \citep{caruana2015intelligible}. A closely related line is \emph{symbolic regression} (SR), which searches over expression trees to recover closed-form models, classically via genetic programming \citep{schmidt09distilling} and can produce highly concise formulas; however, SR is also computationally intensive and scales poorly \citep{ZhangGPU2022}, while KANs, optimized by gradient descent, have been reported to be parameter-efficient and fast to optimize in practice \citep{liu2024kan}.

The need for interpretability is equally pressing in causal inference, where estimators are expected not only to predict counterfactuals but also to justify treatment decisions \citep{athey15recursive}. Early approaches such as causal trees \citep{athey15recursive} provided transparent rule sets, and linear regressors were interpretable through their coefficients as causal parameters \citep{hahn18regularization}. However, partially interpretable algorithms like causal forests \citep{Foster2011SubgroupIF} and BART \citep{hill11bayesian} offered better predictive accuracy, while neural networks advanced performance even further \citep{shalit2017estimating, johansson16learning, schwab2018perfect, yoon2018ganite, yao2018representation} at the cost of opacity. As a result, interpretability has often been relegated to a secondary concern in causal estimation \citep{Rudin2018StopEB}.

Bridging the trade-off between accuracy and transparency remains challenging. Some promising attempts include attention-based transformers highlighting confounder importance \citep{zhang2023exploring}, causal rule learners that yield decision rules but rely on black-box induction \citep{wu2025newcausalrulelearning, bargaglistoffi2024causalruleensembleinterpretable}, fused lasso regression-based methods producing interpretable effect curves \citep{padilla2025causalfusedlassointerpretable}, and, \add{1}{importantly}, NAM-based estimators of average treatment effects \citep{chen22covariate}. \add{1}{} \textit{Model distillation} approaches, such as training a surrogate on top of the work of \citet{shalit2017estimating} \citep{kim2021learninginterpretablemodelscausal}, offer another path but depend on surrogate fidelity. Yet, adressing causal inference and interpretability together requires futher work \citep{moraffah20causal}. In this context, KAN-inspired models such as our approach are promising because they combine the interpretability of additive structures with an expressive power undistinguishable of that of neural networks.

To the best of our knowledge, we are the second work that employs KANs for causal estimation, after KANITE \citep{mehendale2025kanite}, which replaces MLP backbones with KANs to estimate individual treatment effects under multiple treatments, using integral probability metrics/entropy-balancing variants and reporting gains in causal inference metrics over strong baselines. \replace{1}{I}{Although this paper already performs a KAN-ification, it does not provide a systematic method for carrying it out, and i}ts emphasis is predictive accuracy; it neither targets interpretability nor provides closed-form effect functions or visually interpretable plots. Instead, our work focuses on extracting human-readable causal effect formulas.

\add{1}{Last, we also want to compare our proposal with the NAM-based approach of \citet{chen22covariate}. First, while NAMs offer decomposable architectures, they do not provide closed-form symbolic expressions nor a principled pipeline for transforming neural estimators into interpretable representations; second, the method proposed by \citet{chen22covariate} is suitable only for ATE prediction, and third, our pipeline regularizes the functional behavior learned by neural networks, which can exhibit spiking or irregular patterns, while the simplification steps of our pipelie produce smoother, more stable mappings, which directly enhances interpretability and faithfulness in some real world problems \citep{liu2024kan, wang2025kolmogorov}}